% Modernised reading — fonds Alexandre Grothendieck, shelfmark 39 (pages 1-19).
%
% This is an INTERPRETATION, not a transcription. It re-reads the pages in
% today's notation and vocabulary, states the hypotheses the manuscript leaves
% implicit, and drops the critical apparatus so that the mathematics reads
% straight through. Everything the brackets and underlines carried moves into
% a footnote. It is held to being mathematically correct as it stands; where
% the manuscript is loose or mistaken the true statement appears and the
% footnote says what the page has. In any dispute of substance the
% transcription governs: whoever cites Grothendieck cites batch-01.fr.tex.
%
% Scope: a single document for the whole folder. The shelfmark holds one
% batch (pages 1-19), read whole before any of this was written. Pages 1, 8
% and 17 are blank separators.
%
% Notation fixed for the whole folder. k is the base ring (a field in the
% first redaction, pages 2-7; any commutative ring in the second, pages 9-16);
% S_n = k[[t_1,...,t_n]], m_n its ideal (t_1,...,t_n), epsilon : S_n -> k the
% augmentation f |-> f(0), X_n = Spf S_n. The family is written A = (A_n); the
% page writes A in the first redaction and A_* in the second. An
% A-morphism X_m -> X_n is an n-tuple of elements of A_m whose values at 0 are
% nilpotent. The conditions are given NAMES here — (I) coordonnées, (S)
% substitution, (R) anneau, (L) local, (H) inversion/hensélienne, (D)
% coefficients, (E) équations différentielles, (F) feuilletages, (Q) quotients
% — because the page's letters collide: page 11 calls the henselian condition
% d), page 14 opens the coefficient condition also as d), pages 13 and 15
% invoke e), and the first and second redactions use a)-f) for different
% things. A table in the spine section gives the correspondence.
%
% Substantive departures from the page, with pages.
%  - Page 2: the smallest family satisfying stability, subring and t_1 in A_1
%    is Lambda[t], Lambda the prime RING, not P[t] with P the prime field;
%    the field needs the locality condition, which the page struck from that
%    line. Page 2 also writes 1/f in S_n where A_n is meant.
%  - Page 3: the NB « moyennant b) et e), la condition a) équivaut à 1 in A_0,
%    différence, produit » is restated as what holds (given coordinates and
%    substitution, the ring condition is equivalent to those three); the
%    second redaction, page 9, reorganises the axioms exactly this way.
%  - Page 5: the claim that the Hausdorff group law is an A-law is announced
%    and reduced to the linear case by Ado; the linear case is not done, and
%    this reading says what it would still need (the exponential is a
%    time-one map, and time is only a formal variable in (E)).
%  - Page 6: « cos, sin, tg in A_1 » is justified by a route the page does not
%    give (arctan, inverse series, a square root by henselianity).
%  - Page 9: the equivalence « composition closed <=> p^*(A_n) in A_m » needs
%    the constants of A_n in A_m — his own margin note says so.
%  - Page 10: the transversal fibre products announced under a)-c) need the
%    henselian condition (H); « pleinement fidèle (?) » is faithful, not full;
%    Hom_A(X, X_1) is the ideal of functions nilpotent at 0, not all of O_A(X).
%  - Page 11: Cor. 1 as proved needs the substitution t_{n+1} = 1, which (S)
%    does not allow — his own margin objection; the reading proves it with the
%    coefficient condition (D) of page 14, which also gives the page's first
%    form (any unit value at 0) directly, and with it f(0)^{-1} in A_0 — the
%    fact his reductions between the forms use tacitly.
%  - Page 19: the page writes the unit alpha as an element of A_0^*; it is a
%    unit of A_0[[Z]].
%
% Announced and never established in the folder: whether f') implies f)
% (page 3 margin); whether S_n is faithfully flat over A_n (page 12 margin);
% the Hausdorff factorisation (page 5); the bijection between formal quotients
% and formal equivalence relations (page 15, asserted); the condition (Q)
% itself, which is an axiom. The formal preparation theorem of pages 18-19 is
% never connected by the page to the classes A; the connection drawn in the
% résumé and the spine is this reading's.
%
% Pass: Opus 5 (claude-opus-5), 2026-09-13 — first pass, unchecked against the pages by a human.
\documentclass[11pt,a4paper]{article}
% Le préambule commun à toutes les transcriptions du fonds Grothendieck.
%
% Il tient en une idée : **l'incertitude se compose, elle ne se résout pas.**
% Une transcription qui lisse un mot douteux a détruit la seule chose pour
% laquelle on la faisait — un lecteur doit pouvoir savoir, à chaque endroit,
% ce qui était sur le feuillet et ce qui a été ajouté. Les six macros qui
% suivent sont donc l'appareil critique, et non de la décoration.
%
% Les mêmes commandes sont reconnues par `scripts/render.mjs`, qui produit la
% vue de lecture du site. Une macro ajoutée ici doit l'être aussi là-bas,
% faute de quoi le rendu échouera bruyamment — ce qui est voulu : un rendu
% silencieusement faux est pire qu'un rendu absent.

\ProvidesPackage{grothendieck}[2026/08/08 Transcriptions du fonds Grothendieck]

\usepackage{amsmath,amssymb,amsthm}
\usepackage{mathrsfs}
\usepackage{stmaryrd}
% Les diagrammes commutatifs sont partout dans ce fonds : c'est la langue
% même de la géométrie algébrique de Grothendieck, et une transcription qui
% les rendrait en prose ne transcrirait rien.
\usepackage{tikz-cd}
% Français par défaut — c'est la langue des feuillets — mais l'anglais est
% chargé aussi : la traduction et le résumé sont des documents anglais, et la
% typographie française (espaces avant les deux-points) y serait une faute.
% Ces documents basculent par \selectlanguage{english} en tête de corps.
\usepackage[english,french]{babel}
\usepackage{fontspec}
\usepackage{geometry}
\geometry{a4paper,margin=25mm,marginparwidth=28mm}
\usepackage{marginnote}
\usepackage{xcolor}
% ulem, not soul. `soul`'s \st and \ul are built on a character-by-character
% scan that XeLaTeX's font machinery breaks: the document fails with an
% undefined control sequence rather than degrading. ulem does the same two jobs
% and is Unicode-engine safe. `normalem` stops it hijacking \emph, which the
% transcriptions use for Grothendieck's own emphasis.
\usepackage[normalem]{ulem}
\usepackage{hyperref}
\hypersetup{colorlinks=true,linkcolor=black,urlcolor=black,pdfborder={0 0 0}}

% --- Métadonnées du lot -----------------------------------------------------
% Reprises telles quelles par le script de rendu, qui les affiche en tête de
% la vue de lecture. Elles fixent ce que le fichier prétend couvrir : sans
% elles, une transcription flotte sans point d'attache dans un fonds de seize
% mille feuillets.

\newcommand{\folder}[1]{\gdef\grothFolder{#1}}
\newcommand{\batch}[1]{\gdef\grothBatch{#1}}
\newcommand{\leaves}[2]{\gdef\grothFirst{#1}\gdef\grothLast{#2}}
\newcommand{\foldertitle}[1]{\gdef\grothFoldertitle{#1}}
\gdef\grothFolder{?}\gdef\grothBatch{?}\gdef\grothFirst{?}\gdef\grothLast{?}\gdef\grothFoldertitle{}

% --- Le résumé --------------------------------------------------------------
%
% Il ouvre la lecture modernisée, sous le titre « Résumé », et il est composé
% en petit. Ce n'est pas un ornement : c'est le seul passage du document qui
% ne soit pas des mathématiques, et un lecteur doit pouvoir le reconnaître
% avant de l'avoir lu — puis le sauter s'il connaît déjà le sujet, ou s'y
% arrêter s'il ne le connaît pas. Le filet qui le suit marque la reprise du
% corps.

\newenvironment{resume}
  {\small\setlength{\parskip}{0.45em}}
  {\par\medskip\hrule\medskip}

% --- Mots-clés ---------------------------------------------------------------
%
% En anglais, en clôture du résumé de la lecture modernisée : le vocabulaire
% moderne sous lequel chercher ce que les feuillets construisent. Le manifeste
% du site (`npm run manifest`) les extrait pour étiqueter la cote entière —
% cette ligne est la source unique des tags, il n'y a pas de fichier à côté.
\newcommand{\keywords}[1]{\par\smallskip\noindent
  {\footnotesize\textsc{Keywords} --- \textit{#1}}\par}

% --- L'appareil critique ----------------------------------------------------

% Le numéro de feuillet, en marge. C'est l'ancre qui permet de régler un
% désaccord sur une formule en regardant *un* feuillet plutôt qu'une chemise
% de six cents. Le site s'en sert aussi pour tourner les pages du fac-similé
% quand on fait défiler la transcription.
\newcommand{\leaf}[1]{%
  \marginnote{\footnotesize\color{gray!70}\textsf{#1}}%
  \label{leaf:#1}\ignorespaces}

% Illisible : on ne devine pas. Un mot inventé qui se lit comme les autres est
% le pire résultat possible.
\newcommand{\ill}{\textcolor{red!60!black}{[\,\ldots\,]}}

% Lecture douteuse : le mot est donné, et signalé comme incertain.
\newcommand{\uncertain}[1]{\uline{#1}}

% Ajout de l'éditeur — une lettre manquante, un mot sous-entendu.
\newcommand{\add}[1]{\textcolor{blue!50!black}{[#1]}}

% Biffé par Grothendieck lui-même. Ce qu'il a rayé fait partie du document :
% une rature dit souvent où la pensée a changé de direction.
\newcommand{\struck}[1]{\sout{#1}}

% Note du transcripteur, sur l'état matériel du feuillet.
\newcommand{\note}[1]{\par\smallskip{\small\color{gray!60!black}\textit{[#1]}}\par\smallskip}

% Note marginale de Grothendieck, distincte du corps du texte.
\newcommand{\marginal}[1]{\par\smallskip\noindent\hspace{1em}%
  {\small\color{gray!60!black}#1}\par\smallskip}

% --- Titre ------------------------------------------------------------------

\AtBeginDocument{%
  \begin{center}
    {\large\bfseries Fonds Alexandre Grothendieck --- cote n\textsuperscript{o}~\grothFolder}\\[2pt]
    {\itshape\grothFoldertitle}\\[4pt]
    {\small feuillets \grothFirst--\grothLast\ (lot \grothBatch)}\\[2pt]
    {\footnotesize\color{gray!60!black}Transcription établie d'après le fac-similé numérisé
     par l'Université de Montpellier.\\
     Les lectures incertaines sont soulignées, les illisibles notées [\,\ldots\,],
     les ajouts entre crochets.}
  \end{center}
  \vspace{1em}}

\endinput


\folder{39}
\pages{1}{19}
\dating{[vers 1960-1970]}
\watermark{Édition de démonstration\\interprétation personnelle de l'œuvre}
\foldertitle{Fonctions et structures « élémentaires » [Schémas formels, théorème de préparation] : notes manuscrites (s.d.) — lecture modernisée du dossier entier}

\begin{document}

\section*{Résumé}

\begin{resume}
Qu'est-ce qu'une classe « raisonnable » de fonctions ? Les polynômes, les
fractions rationnelles, les fonctions algébriques, les fonctions qu'on obtient
en résolvant des équations différentielles : chacune de ces familles est plus
riche que la précédente, et chacune est stable par un certain nombre
d'opérations — additionner, multiplier, composer, inverser une application,
parfois intégrer. Ces dix-neuf pages cherchent à dire, une fois pour toutes,
quelles opérations une telle famille doit supporter pour qu'on puisse faire de
la géométrie avec elle.

Le cadre est celui des séries formelles : une « fonction » de $n$ variables
est une série entière $f(t_1, \ldots, t_n)$ dont on ne demande pas qu'elle
converge. Toutes les séries formelles forment un monde trop grand et sans
contenu propre ; une classe de fonctions sera un choix, pour chaque $n$, d'un
sous-anneau $A_n$ de séries. Grothendieck dresse la liste des exigences :
contenir les coordonnées, être stable par substitution, puis par inversion des
changements de variables — et découvre qu'une telle liste, si on la prend au
sérieux, force la famille à ressembler aux fonctions algébriques. Il en tire
la conséquence qui compte : on peut alors définir des « variétés » dont les
changements de cartes sont dans la classe, comme on définit les variétés
analytiques en exigeant des changements de cartes analytiques plutôt que
seulement différentiables. Il pousse ensuite la liste jusqu'aux sous-variétés,
aux quotients, au calcul différentiel, et — en caractéristique nulle — aux
exponentielles de champs de vecteurs, aux groupes de Lie formels et aux
feuilletages.

Le dossier contient deux rédactions de ce programme, la première sur un corps
(pages 2 à 7), la seconde sur un anneau quelconque (pages 9 à 16), plus
ordonnée et plus prudente : on y voit Grothendieck s'objecter en marge à un
pas de sa propre démonstration, et réorganiser ses axiomes pour que chacun ne
serve qu'une fois. Les deux dernières pages démontrent le théorème de
préparation de Weierstrass sous forme formelle : toute série se divise par une
série dont les coefficients d'indice $< N$ sont « petits » et celui d'indice
$N$ inversible, avec un reste polynomial de degré $< N$. Les feuillets ne disent pas pourquoi ce théorème
est là ; mais c'est exactement l'outil qui, dans la théorie qu'on a
développée depuis, sert d'axiome central à ces classes de fonctions.

Car la question a eu une postérité. Les familles de séries stables par
substitution, par inversion et par division de Weierstrass sont aujourd'hui
appelées \emph{systèmes de Weierstrass} ; elles servent en théorie des modèles
et en géométrie modérée, et les séries formelles algébriques, qui sont la plus
petite famille qu'impose la liste de ces pages, sont l'objet du théorème
d'approximation d'Artin. L'autre fil — des changements de cartes pris dans une
classe donnée — est celui des structures définies par un pseudogroupe.

\keywords{Weierstrass system, Weierstrass division theorem, henselian pair, algebraic power series, pseudogroup structure, formal Frobenius theorem, Baker–Campbell–Hausdorff formula, Hasse derivative}
\end{resume}

\pagerange{2}{19}
\section*{Le fil du dossier, et les conventions}

Le dossier ne porte aucun titre de sa main avant la page 18 ; les pages 1, 8
et 17 sont des feuillets intercalaires vierges, qui séparent trois ensembles.

\begin{itemize}
  \item \textbf{Pages 2 à 7 : première rédaction, sur un corps $k$.} Les
    conditions algébriques et leurs plus petits modèles (pages 2 et 3) ; la
    condition d'inversion et l'hensélianité (page 3) ; en caractéristique
    nulle, exponentielles, équations différentielles et groupes formels
    (pages 4 et 5) ; primitives, fonctions usuelles et feuilletages (page 6),
    dont la page 7 est une reprise barrée.
  \item \textbf{Pages 9 à 16 : seconde rédaction, sur un anneau commutatif
    $k$.} Les conditions élémentaires (pages 9 et 10) ; la condition
    hensélienne et ses trois corollaires (pages 11 et 12) ; sous-schémas
    formels (page 13) ; la condition sur les coefficients et le calcul
    différentiel (page 14) ; quotients et condition intégrale (pages 15 et
    16).
  \item \textbf{Pages 18 et 19 : le théorème de préparation formel}, énoncé
    et démontré.
\end{itemize}

La seconde rédaction n'est pas une copie de la première : elle change de
base, et surtout elle remet les axiomes en ordre. On les garde distincts
ci-dessous, et l'on désigne les conditions par des noms plutôt que par les
lettres de la page, qui se chevauchent\footnote{La page 11 appelle d) la
condition hensélienne ; la page 14 ouvre la condition sur les coefficients
sous la lettre d) elle aussi ; les pages 13 et 15 invoquent une condition e) ;
et les lettres a) à f) ne désignent pas les mêmes énoncés dans les deux
rédactions.}. La correspondance est la suivante.

\begin{itemize}
  \item $(\mathrm{I})$ \emph{coordonnées} : $t_i \in A_n$ pour $1 \leq i \leq n$ —
    page 2, d) avec a) ; page 9, a).
  \item $(\mathrm{S})$ \emph{substitution} (clôture par composition) — page 3,
    e) ; page 9, b).
  \item $(\mathrm{R})$ \emph{anneau} : $A_n$ sous-anneau unitaire de $S_n$ —
    page 2, b) ; page 9, c).
  \item $(\mathrm{L})$ \emph{localité} — page 2, c) ; dans la seconde
    rédaction, conséquence de $(\mathrm{H})$ (page 11, Cor.~1), avec la condition
    sur les coefficients.
  \item $(\mathrm{H})$ \emph{inversion}, ou condition hensélienne — page 3,
    f) ; page 11, d).
  \item $(\mathrm{D})$ \emph{coefficients} — page 14, « d) ».
  \item $(\mathrm{E})$ \emph{équations différentielles}, en caractéristique
    nulle — page 4, g) et g$'$).
  \item $(\mathrm{F})$ \emph{feuilletages} — page 6, h) ; et sa forme
    $(\mathrm{Q})$ \emph{quotients}, ou condition intégrale — page 15, f).
\end{itemize}

\textbf{Conventions.} $k$ est l'anneau de base ; $S_n = k[[t_1, \ldots, t_n]]$,
$\mathfrak{m}_n = (t_1, \ldots, t_n)$, $\varepsilon : S_n \to k$ l'augmentation
$f \mapsto f(0)$, et $X_n = \mathrm{Spf}\, S_n$, le germe formel de l'espace
affine de dimension $n$ à l'origine. Une classe est une famille
$A = (A_n)_{n \geq 0}$ de parties $A_n \subset S_n$\footnote{La page écrit $A$
dans la première rédaction, $A_*$ dans la seconde.}. Un morphisme de
$k$-schémas formels $X_m \to X_n$ est un homomorphisme continu
$S_n \to S_m$, c'est-à-dire un $n$-uplet $(f_1, \ldots, f_n)$ d'éléments de
$S_m$ dont les valeurs $f_i(0)$ sont nilpotentes dans $k$ ; sur un corps,
nulles\footnote{La page 2 dit seulement « donné par $f_1, \ldots, f_n \in S_m$ » ;
la condition sur $\varepsilon(f_i)$ est écrite page 9. Elle est nécessaire : un
homomorphisme $S_n \to S_m$, $t_i \mapsto f_i$, n'est continu, et une
substitution $\varphi(f_1, \ldots, f_n)$ n'a de sens, que si les $f_i$ sont
topologiquement nilpotents.}. C'est un \emph{$A$-morphisme} si les $f_i$ sont
dans $A_m$.

\textbf{Le point de vue.} Les objets $X_n$ et les $A$-morphismes forment,
sous les conditions ci-dessous, une catégorie à produits finis
($X_m \times X_n = X_{m+n}$) dont les automorphismes sont les « changements de
coordonnées admissibles » : ce que Grothendieck appelle un \emph{$A$-schéma
formel} est un germe formel lisse muni d'une classe de systèmes de coordonnées
définis à un tel changement près — une structure de pseudogroupe, au sens
d'Ehresmann, dans le cadre formel\footnote{Le mot « pseudogroupe » n'est pas
sur la page. Les structures définies par un pseudogroupe de transformations
(Lie, Cartan, Ehresmann dans les années 1950) lui étaient connues ;
l'analogie est de cette lecture.}.

\textbf{Ce que le dossier annonce sans l'établir.} Que la condition
d'hensélianité entraîne celle d'inversion (page 3) ; que $S_n$ soit fidèlement
plat sur $A_n$ (page 12) ; que la loi de Hausdorff soit une $A$-loi (page 5) ;
la correspondance entre quotients formels et relations d'équivalence
formelles (page 15). Et surtout le lien entre le théorème de préparation des
pages 18 et 19 et les classes $A$ : les feuillets ne le font pas.

\pagerange{2}{3}
\section*{Première rédaction : les conditions algébriques sur un corps (pages 2 et 3)}

Soit $k$ un corps. Les anneaux $S_n$ sont reliés entre eux par les
substitutions de variables : toute application $\sigma : \{1, \ldots, n\} \to
\{1, \ldots, m\}$ définit $S_n \to S_m$, $t_i \mapsto t_{\sigma(i)}$. Les $S_n$
forment ainsi un foncteur de la catégorie des ensembles finis (toutes
applications, ensemble vide compris) vers les anneaux\footnote{La page parle
d'« anneau cosimplicial ». Les conditions qu'elle impose — permutations,
ajout d'une variable, identification de deux variables — engendrent toutes
les applications entre ensembles finis, pas seulement les applications
croissantes : c'est plus qu'une structure cosimpliciale. La page 9 ajoute
« coaugmenté », qui correspond à l'ensemble vide, $S_0 = k$.}. On se donne
$A_n \subset S_n$, et l'on demande :

\begin{itemize}
  \item la \emph{stabilité} : $A = (A_n)$ est un sous-foncteur — $A_n$ est
    stable par $\mathfrak{S}_n$, $f(t_1, \ldots, t_n) \in A_n$ entraîne
    $f \in A_{n+1}$ (une variable $t_{n+1}$ qui ne figure pas), et
    $f(t_1, \ldots, t_{n-1}, t_{n-1}) \in A_{n-1}$ (deux variables identifiées) ;
  \item $(\mathrm{R})$ : chaque $A_n$ est un sous-anneau unitaire de $S_n$ ;
  \item $(\mathrm{L})$ : si $f \in A_n$ et $f(0) \neq 0$, alors $f^{-1} \in A_n$
    \footnote{La page écrit « alors $\frac{1}{f} \in S_n$ », ce qui est toujours vrai
    ; c'est $A_n$ qui est voulu, comme le montre l'équivalence qu'elle écrit
    ensuite.} ;
    de façon équivalente, $A_n$ est un anneau local et l'inclusion
    $A_n \to S_n$ est un homomorphisme local ;
  \item $(\mathrm{I})$ : $t_1 \in A_1$ (et donc, par stabilité, $t_i \in A_n$).
\end{itemize}

L'équivalence dans $(\mathrm{L})$ est immédiate : si tout élément de $A_n$ hors
de $\mathfrak{m}_n$ est inversible dans $A_n$, alors $A_n \cap \mathfrak{m}_n$ est
un idéal dont le complémentaire est formé d'unités, donc l'unique idéal
maximal ; réciproquement, si $A_n$ est local et l'inclusion locale, un élément
hors de $\mathfrak{m}_n$ n'est pas dans l'idéal maximal de $A_n$.

\textbf{Les plus petits modèles.} Notons $\Lambda$ l'image de $\mathbf{Z}$ dans
$k$ ($\mathbf{Z}$ ou $\mathbf{F}_p$) et $P$ le sous-corps premier de $k$
($\mathbf{Q}$ ou $\mathbf{F}_p$).

\begin{itemize}
  \item Stabilité et $(\mathrm{R})$ : $A_n = \Lambda$, les constantes.
  \item Stabilité, $(\mathrm{R})$, $(\mathrm{L})$ : $A_n = P$.
  \item Stabilité, $(\mathrm{R})$, $(\mathrm{I})$ : $A_n = \Lambda[t_1, \ldots, t_n]$
    \footnote{La page écrit ici $P[t_1, \ldots, t_n]$, après avoir biffé
    c) de la liste des conditions. Sans la condition de localité, rien
    n'oblige à inverser les entiers : en caractéristique nulle, la plus petite
    famille est $\mathbf{Z}[t_1, \ldots, t_n]$, et $\mathbf{Q}[t_1, \ldots, t_n]$
    n'est pas minimale.}.
  \item Les quatre : $A_n = P[t_1, \ldots, t_n]_{\mathfrak{p}_n}$, localisé en
    $\mathfrak{p}_n = (t_1, \ldots, t_n)$ — les fractions rationnelles à
    coefficients dans $P$ définies à l'origine.
\end{itemize}

\textbf{La composition.} La condition suivante porte sur les morphismes :

\begin{itemize}
  \item $(\mathrm{S})$ : le composé de deux $A$-morphismes
    $X_m \to X_n \to X_p$ est un $A$-morphisme.
\end{itemize}

Moyennant $(\mathrm{S})$, on a
$\mathrm{Hom}_A(X_m, X_n) \simeq \mathrm{Hom}(A_n, A_m)$, les homomorphismes
locaux de $A_0$-algèbres\footnote{C'est la note marginale de la page 3,
« $\mathrm{Hom}_A(X_m, X_n) \simeq \mathrm{Hom}_{A_0\text{-alg}}(A_n, A_m)$ »,
dont le début est d'une lecture incertaine. L'isomorphisme demande de se
restreindre aux homomorphismes qui envoient les $t_i$ dans l'idéal
$A_m \cap \mathfrak{m}_m$ et se prolongent continûment aux $S$ ; nous l'énonçons
ainsi.}. Les deux familles $P[t]$ et $P[t]_{\mathfrak{p}}$ la vérifient :
substituer des polynômes, ou des fractions régulières nulles à l'origine,
dans une fraction régulière à l'origine en donne une.

Grothendieck observe alors qu'une partie des axiomes peut être remplacée par
une condition sur trois morphismes seulement\footnote{La page écrit :
« Moyennant b) et e), la condition a) équivaut au fait que $1 \in A_0$, le
morphisme “différence” $X_1 \times X_1 \to X_1$ est un $A$-morphisme, le
morphisme “produit” aussi. » Telle quelle, l'équivalence ne tient pas : ces
trois conditions disent que les $A_n$ sont des anneaux, pas qu'ils sont
stables par les substitutions de variables. L'énoncé correct est celui que
nous donnons, et c'est exactement la réorganisation que fait la seconde
rédaction, page 9.}. Si les coordonnées sont dans $A$ et que $A$ est stable
par substitution au sens fort — $\varphi \in A_n$ et $f_1, \ldots, f_n \in A_m$
nuls à l'origine entraînent $\varphi(f_1, \ldots, f_n) \in A_m$ —, la stabilité
est automatique (les substitutions de variables sont des substitutions par
des coordonnées), et $(\mathrm{R})$ équivaut à
\[
  1 \in A_0, \qquad t_1 - t_2 \in A_2, \qquad t_1 t_2 \in A_2,
\]
c'est-à-dire : la différence et le produit $X_1 \times X_1 \to X_1$ sont des
$A$-morphismes. Autrement dit, $X_1$ est un objet anneau de la catégorie des
$A$-morphismes.

\pagerange{3}{3}
\section*{Inversion et hensélianité (page 3)}

La condition qui donne son contenu à la théorie est la suivante.

\begin{itemize}
  \item $(\mathrm{H})$ : si $u : X_n \to X_n$ est un $A$-morphisme et un
    isomorphisme — c'est-à-dire si le jacobien $\det(\partial u_i/\partial t_j)(0)$
    est non nul —, alors $u^{-1}$ est un $A$-morphisme.
\end{itemize}

C'est l'analogue, pour les morphismes, de la condition $(\mathrm{L})$ pour les
fonctions : $(\mathrm{L})$ inverse une fonction, $(\mathrm{H})$ inverse un
changement de coordonnées. Elle entraîne :

\begin{itemize}
  \item $(\mathrm{H}')$ : les anneaux locaux $A_n$ sont henséliens.
\end{itemize}

Plus précisément, $(\mathrm{H})$ en dimension $n$ entraîne que $A_{n-1}$ est
hensélien, par l'argument que la seconde rédaction détaille (page 12,
corollaire 3) : une équation $F(t, Z) = 0$ à racine simple se résout en
inversant le changement de coordonnées $(t, z) \mapsto (t, F(t, z))$. La
réciproque — $(\mathrm{H}')$ entraîne-t-il $(\mathrm{H})$ ? — est posée en marge
et n'est pas tranchée par les feuillets\footnote{La marge ajoute « il suffit
par récurrence », d'une lecture incertaine, sans suite.}. Un passage barré
esquissait une description de $u^{-1}$ par une clôture hensélienne.

\textbf{Le plus petit modèle hensélien.} La plus petite famille satisfaisant
stabilité, $(\mathrm{R})$, $(\mathrm{L})$, $(\mathrm{I})$ et $(\mathrm{H}')$ est
\[
  A_n = P\{t_1, \ldots, t_n\},
\]
l'hensélisé de $P[t_1, \ldots, t_n]_{\mathfrak{p}_n}$, et elle satisfait
$(\mathrm{S})$ et $(\mathrm{H})$. L'hensélisé se plonge dans $S_n$ : c'est
l'anneau des séries formelles à coefficients dans $P$ qui sont algébriques sur
$P[t_1, \ldots, t_n]$, et la minimalité découle de sa propriété universelle
(un homomorphisme local vers un anneau hensélien se prolonge de façon unique à
l'hensélisé)\footnote{La page note $P\{t_1, \ldots, t_n\}$ et l'identifie à
la clôture hensélienne ; l'identification aux séries algébriques, et le nom de
séries formelles algébriques, sont de cette lecture. C'est le cadre du
théorème d'approximation d'Artin (1969).}.

Moyennant ces conditions, conclut la page, on peut définir la catégorie des
$A$-schémas formels, les Modules sur eux, et le calcul différentiel
(opérateurs différentiels, parties principales)\footnote{Pour les parties
principales, il faut davantage : c'est la condition sur les coefficients que
la seconde rédaction introduit page 14, précisément dans ce but. Une note
marginale en bas de page mentionne des produits fibrés « transversaux » ; elle
est en grande partie illisible.}.

\pagerange{4}{5}
\section*{En caractéristique nulle : exponentielles et groupes formels (pages 4 et 5)}

On suppose désormais $k$ de caractéristique nulle. Un champ de vecteurs formel
$\xi$ sur un schéma formel lisse $X$ a alors une exponentielle,
\[
  X_1 \times X \longrightarrow X, \qquad (t, x) \longmapsto \exp(t\xi) \cdot x
  = \sum_{r \geq 0} \frac{t^r}{r!}\, \xi^r(x),
\]
où $t$ est le « temps », une variable formelle. On demande :

\begin{itemize}
  \item $(\mathrm{E})$ : si $X$ est un $A$-schéma formel et $\xi$ un $A$-champ
    de vecteurs, le flot $X_1 \times X \to X$ est un $A$-morphisme ; et plus
    généralement, pour un champ dépendant du temps
    $\xi : X_1 \times X \to \mathbf{T}(X)$, la solution $x = \varphi(t, x_0)$ de
    l'équation différentielle formelle
    \[
      \frac{dx}{dt} = \xi(t, x), \qquad x(0) = x_0,
    \]
    est un $A$-morphisme $X_1 \times X \to X$.
\end{itemize}

Il suffit de le demander pour $X = X_n$. La seconde forme, remarque la page,
fait du temps une variable explicite ; c'est elle qui sert.

\textbf{Groupes de Lie formels.} Soit $k_0 = A_0$ ; par $(\mathrm{L})$, c'est un
sous-corps de $k$. Pour une algèbre de Lie $\mathfrak{g}_0$ de dimension finie
sur $k_0$, la formule de Baker--Campbell--Hausdorff définit une loi de groupe
formel sur le complété à l'origine $\widehat{W}(\mathfrak{g}_0)$ de l'espace
vectoriel $\mathfrak{g}_0 \otimes_{k_0} k$ — un $A$-schéma formel, une base de
$\mathfrak{g}_0$ fournissant des coordonnées. On obtient un foncteur
\[
  \mathcal{L}(k_0) \longrightarrow \{\text{groupes formels sur } k_0\}
  \longrightarrow \{\text{groupes formels sur } k\},
\]
la seconde flèche étant l'extension des scalaires. Grothendieck annonce que ce
foncteur se factorise par les $A$-groupes formels : la loi de Hausdorff est
une $A$-loi, et un homomorphisme $\mathfrak{g}_0 \to \mathfrak{h}_0$ donne un
$A$-homomorphisme. Par le théorème d'Ado, il ramène l'énoncé au cas d'une
sous-algèbre $\mathfrak{g}_0 \subset \mathfrak{gl}(n, k_0)$ : il suffirait de
savoir que le sous-groupe formel de $\mathrm{GL}_n$ qu'elle définit est un
sous-$A$-schéma formel.

La page s'arrête là, sur des points de suspension, et le cas linéaire n'est pas
fait\footnote{Deux passages, barrés, esquissaient un argument par le système
de Pfaff associé ; ils sont illisibles pour l'essentiel. Ce qui reste à faire
n'est pas immédiat. En coordonnées de seconde espèce,
$(s_1, \ldots, s_d) \mapsto \exp(s_1 e_1) \cdots \exp(s_d e_d)$, chaque facteur
est le flot d'un champ linéaire au temps $s_i$ et relève donc de
$(\mathrm{E})$ ; mais la loi de Hausdorff est écrite en coordonnées canoniques,
et passer de l'une à l'autre demande que $x \mapsto \exp(x)$ soit une
$A$-application — le flot \emph{au temps $1$}, que $(\mathrm{E})$ ne fournit
pas, le temps n'y étant qu'une variable formelle qu'on ne peut spécialiser
en $1$.}. Nous le laissons pour ce qu'il est : un énoncé annoncé.

\pagerange{6}{7}
\section*{Primitives, fonctions usuelles, feuilletages (pages 6 et 7)}

\textbf{Corollaire.} $A_1$ est stable par primitivation nulle à l'origine, et
contient $\log(1 + t)$, $\exp t$, $\cos t$, $\sin t$, $\tan t$.

La primitive $\int_0^t f$ d'un $f \in A_1$ est la solution, pour
$x_0 = 0$, de l'équation $dx/dt = f(t)$ sur $X_1$, dont le second membre est
dans $A_2$ ; elle est dans $A_1$ par $(\mathrm{E})$. Ensuite\footnote{La page
donne $\log(1+t)$ comme primitive de $\frac{1}{1+t}$ et $\exp$ comme « série
inverse », mais ne dit pas comment on obtient $\cos$, $\sin$ et $\tan$ ; le
chemin qui suit est de cette lecture.} :

\begin{itemize}
  \item $\frac{1}{1+t} \in A_1$ par $(\mathrm{L})$, donc $\log(1+t) \in A_1$ ; sa
    série réciproque $e^t - 1$ est dans $A_1$ par $(\mathrm{H})$, donc $e^t$ aussi ;
  \item de même $\arctan t = \int_0^t \frac{ds}{1+s^2} \in A_1$, et sa série
    réciproque $\tan t$ est dans $A_1$ ;
  \item $\cos t = (1 + \tan^2 t)^{-1/2}$ est la racine de $Z^2 (1 + \tan^2 t) = 1$
    qui vaut $1$ en $0$, racine simple : elle est dans $A_1$ par
    hensélianité ; et $\sin t = \tan t \cdot \cos t$.
\end{itemize}

\textbf{Feuilletages.} Soit $X$ un schéma formel lisse de dimension $n = p + q$.
Un feuilletage de dimension $p$ de $X$ est la donnée d'un quotient formel lisse
$X \to B$, submersif, avec $\dim B = q$. Il lui correspond le sous-module
$\mathfrak{P} \subset \Theta(X)$ des champs de vecteurs tangents aux fibres, qui
est un sous-fibré de rang $p$ involutif, et en caractéristique nulle, par le
théorème de Frobenius formel, tout sous-fibré involutif s'obtient ainsi\footnote{La
page dit « système de Pfaff de dimension $p$ complètement intégrable ». Un
système de Pfaff désigne classiquement un système de $1$-formes ; ici il
s'agit de champs de vecteurs, ce qu'on appelle aujourd'hui une distribution
involutive. L'hypothèse de caractéristique nulle, posée page 4, est
indispensable : en caractéristique $p$ le théorème de Frobenius formel est
faux.}.

Si $X$ est un $A$-schéma formel, un feuilletage est un \emph{$A$-feuilletage}
s'il existe sur $B$ une $A$-structure faisant de $X \to B$ un $A$-morphisme ;
elle est alors unique, par $(\mathrm{H})$. Le sous-fibré $\mathfrak{P}$ d'un
$A$-feuilletage est engendré par des $A$-champs. On demande la réciproque :

\begin{itemize}
  \item $(\mathrm{F})$ : un sous-fibré involutif engendré par des $A$-champs
    définit un $A$-feuilletage. De façon équivalente : si $i : Y \to X$ est un
    $A$-morphisme transverse à $\mathfrak{P}$ en $0$ — le composé
    $\mathbf{T}_0(Y) \to \mathbf{T}_0(X) \to \mathbf{T}_0(X)/\mathfrak{P}(0)$ est un
    isomorphisme —, l'unique rétraction $\pi : X \to Y$ constante sur les
    feuilles est un $A$-morphisme.
\end{itemize}

Il en résulte que la feuille passant par $0$ est un sous-$A$-schéma formel.
$(\mathrm{F})$ entraîne $(\mathrm{E})$, et pour $p = 1$ lui est
équivalente\footnote{C'est la note marginale « h) implique g), qui équivaut
à h) pour $p = 1$ ». Pour $p = 1$ et un champ non nul à l'origine, redresser
le champ revient à connaître son flot, et $(\mathrm{H})$ fait passer de l'un à
l'autre ; un champ quelconque se ramène à ce cas en lui ajoutant
$\partial/\partial s$ sur $X_1 \times X$. La page ne détaille pas.}.

La page 7 reprend la fin de la page 6 et a été barrée tout entière. Elle y
formule l'équivalence de deux conditions pour une rétraction $\pi$ comme
ci-dessus : $\pi$ est un $A$-morphisme ; il existe un $A$-isomorphisme
$X \simeq B \times Y$ qui transforme le feuilletage en celui de la première
projection. Une parenthèse, encadrée et biffée, construisait d'abord un
isomorphisme formel $X \simeq Y \times Z$ à partir de deux transversales.

\pagerange{9}{10}
\section*{Seconde rédaction : les conditions élémentaires sur un anneau (pages 9 et 10)}

La seconde rédaction reprend tout sur un anneau commutatif $k$ quelconque. Les
$X_n$ forment un schéma formel simplicial augmenté, les $S_n$ un anneau
topologique cosimplicial coaugmenté, et l'on pose les trois conditions
élémentaires :

\begin{itemize}
  \item $(\mathrm{I})$ \emph{identités} : les opérations simpliciales entre les
    $X_n$ sont des $A$-morphismes ; de façon équivalente, $t_i \in A_n$ pour
    $1 \leq i \leq n$, ou encore les identités des $X_n$ sont des $A$-morphismes ;
  \item $(\mathrm{S})$ \emph{composition} : si $\varphi \in A_n$ et si
    $f_1, \ldots, f_n \in A_m$ ont des valeurs $f_i(0)$ nilpotentes, alors
    $\varphi(f_1, \ldots, f_n) \in A_m$ ;
  \item $(\mathrm{R})$ \emph{anneau} : $1 \in A_0$, et la différence et le
    produit $X_1 \times X_1 \to X_1$ sont des $A$-morphismes ; c'est dire
    $t_1 - t_2,\ t_1 t_2 \in A_2$, ou encore que les $A_n$ sont des
    sous-anneaux unitaires des $S_n$.
\end{itemize}

La page donne $(\mathrm{S})$ sous la forme « le composé de deux $A$-morphismes
est un $A$-morphisme », qu'elle déclare équivalente à la forme ci-dessus. Les
deux ne coïncident que si les constantes suivent : la clôture par composition
ne porte que sur des fonctions nulles à l'origine, et pour l'étendre à un
$\varphi$ quelconque il faut que $\varphi(0)$ soit dans $A_m$, ce qui est vrai dès
que $\mathrm{Im}(A_0 \to S_m) \subset A_m$\footnote{C'est exactement la note
marginale de la page 9 : « ce serait vrai (sans supposer $\varepsilon(\varphi)$
nilpotent) \ldots\ sachant que $\mathrm{Im}(A_0 \to S_n) \subset A_n$ ». Une
autre note, en haut, porte sur les valeurs $\varepsilon(f_i)$ et se lit mal.
La condition $(\mathrm{R})$ est en outre encadrée et barrée sur la page, sans
rien qui la remplace ; la page 10 la tient pour acquise, et nous aussi.}. Nous
prenons $(\mathrm{S})$ sous la forme forte.

Ces trois conditions font de $A$ un sous-anneau cosimplicial de $S$, stable par
toutes les substitutions de variables. La plus petite classe est
\[
  A_n = \mathrm{Im}\bigl(\mathbf{Z}[t_1, \ldots, t_n] \to k[[t_1, \ldots, t_n]]\bigr)
  = \Lambda[t_1, \ldots, t_n],
\]
où $\Lambda$ est l'image de $\mathbf{Z}$ dans $k$, isomorphe à $\mathbf{Z}$ ou à
$\mathbf{Z}/c\mathbf{Z}$, $c$ la caractéristique de $k$.

\textbf{La catégorie des $A$-schémas formels.} Les $X_n$ et les
$A$-morphismes forment une catégorie $\mathrm{Sfl}_A$, avec produits finis
$X_n \simeq X_1^n$. La page y annonce aussi des produits fibrés dans le cas
transversal\footnote{Pour paramétrer un produit fibré transversal par des
$A$-coordonnées, il faut le théorème des fonctions implicites dans la classe,
c'est-à-dire la condition hensélienne $(\mathrm{H})$ de la page 11, que la
page n'a pas encore introduite.}.

Soit $k_0 = A_0$. On a un foncteur $\mathbb{V}_A$ de la catégorie des
$k_0$-modules libres de type fini vers les $A$-schémas formels,
$M_0 \mapsto \mathrm{Spf}\, \widehat{\mathrm{Sym}}_k(M_0 \otimes_{k_0} k)$, qui
est contravariant\footnote{La page dit « un foncteur » sans préciser la
variance ; $M \mapsto \mathrm{Spf}\,\widehat{\mathrm{Sym}}(M)$ renverse le sens des
flèches.}, fidèle, et qui transforme les objets en groupe en $A$-groupes
formels\footnote{La page écrit « pleinement fidèle », suivi d'un point
d'interrogation. Il n'est pas plein : $t \mapsto t^2$ est un $A$-morphisme
$X_1 \to X_1$ qui ne vient d'aucune application linéaire.}. Plus
généralement, le complété à l'origine d'un $k_0$-schéma lisse pointé porte une
$A$-structure, et ce foncteur commute aux produits finis, donc envoie groupes
sur groupes\footnote{Hors du cas des espaces affines, les changements de
coordonnées étales qui interviennent sont des séries algébriques, et c'est
encore la condition hensélienne qui les met dans $A$ ; la page 13 fait la
construction sous cette hypothèse.}.

La marge de la page 10 note l'anneau $\mathcal{O}_A(X) \subset \mathcal{O}(X)$ des
$A$-fonctions sur $X$, et l'identifie à $\mathrm{Hom}_A(X, X_1)$ au sein de
$\mathrm{Hom}(X, X_1)$. Il faut distinguer : $\mathrm{Hom}(X, X_1)$ n'est que
l'idéal des fonctions nilpotentes à l'origine, et c'est
$\mathrm{Hom}_A(X, X_1) = \mathcal{O}_A(X) \cap \mathrm{Hom}(X, X_1)$ qui est vrai ;
$\mathcal{O}_A(X)$ lui-même est l'anneau qui vaut $A_n$ en coordonnées. La même
marge mentionne les Modules $A$-formels et les fibrés vectoriels $A$-formels,
et affirme que le foncteur « anneau des $A$-fonctions » commute aux limites
projectives, sans préciser lesquelles.

\pagerange{11}{12}
\section*{La condition hensélienne et ses corollaires (pages 11 et 12)}

\begin{itemize}
  \item $(\mathrm{H})$ \emph{condition hensélienne} : si $u : X_n \to X_n$ est un
    $A$-morphisme qui est un isomorphisme formel — si
    $\det\bigl(\frac{\partial f_i}{\partial t_j}(0)\bigr)$ est inversible dans $k$ —,
    alors $u^{-1}$ est un $A$-morphisme. Autrement dit, le foncteur d'oubli des
    $A$-schémas formels vers les schémas formels sur $k$ est conservatif.
\end{itemize}

Sur un anneau, un endomorphisme $(f_1, \ldots, f_n)$ de $X_n$ dont les $f_i(0)$
sont nilpotents et dont le jacobien à l'origine est inversible est bien un
automorphisme : on écrit $f = c + h$ avec $c = f(0)$ nilpotent et $h(0) = 0$ ;
$h$ est inversible par le théorème d'inversion formel, et $h^{-1}(x - c)$
converge.

\textbf{Corollaire 1.} Si $f \in A_n$ et si $f(0)$ est inversible dans $k$,
alors $f^{-1} \in A_n$. En particulier $A_n \cap \mathfrak{m}_n$ est contenu dans
le radical de Jacobson de $A_n$ (prendre $f = 1 + ag$, $g \in A_n \cap
\mathfrak{m}_n$).

La démonstration de la page considère l'automorphisme
$u(x, s) = (x, f(x)\, s)$ de $X_{n+1}$, de jacobien $f(0)$ à l'origine, d'inverse
$(x, s) \mapsto (x, f(x)^{-1} s)$. Par $(\mathrm{H})$,
$g = f^{-1}\, t_{n+1} \in A_{n+1}$. La page substitue alors $t_{n+1} = 1$ pour
conclure — et objecte elle-même, en marge, que $1$ n'est pas nilpotent, de
sorte que $(\mathrm{S})$ n'autorise pas cette substitution ; elle propose
d'exiger une propriété de divisibilité, $t_{n+1} f \in A_{n+1} \Rightarrow
f \in A_n$\footnote{Une autre note marginale, en face de l'énoncé, parle d'une
« condition de coefficientabilité » ; la lecture en est incertaine.}. La
condition sur les coefficients $(\mathrm{D})$ de la page 14 suffit : $f^{-1}$ est
le coefficient de $t_{n+1}$ dans $g$, donc est dans $A_n$. C'est sous cette
hypothèse supplémentaire que nous tenons le corollaire pour démontré.

La page passe ensuite, par des équivalences, aux formes « $f(0) = 1$ » et
« $A_n \cap \mathfrak{m}_n \subset \mathrm{rad}\, A_n$ ». Revenir de là à la
première forme suppose $f(0)^{-1} \in A_0$ ; c'est vrai, mais c'est le cas
$n = 0$ de la première forme, donc de l'argument ci-dessus, et non une
conséquence des seules formes affaiblies.

\textbf{Corollaire 2.} Soit $P(Z) = Z^N + f_1 Z^{N-1} + \cdots + f_N \in A_n[Z]$
unitaire, et $\zeta \in k_0$ une racine du polynôme réduit
$P^\varepsilon(Z) = Z^N + a_1 Z^{N-1} + \cdots + a_N$, $a_i = f_i(0) \in k_0$, telle
que $(P^\varepsilon)'(\zeta)$ soit inversible dans $k$. L'unique $F \in S_n$ tel que
$P(F) = 0$ et $F(0) = \zeta$ est dans $A_n$. En d'autres termes, le couple
$(A_n, A_n \cap \mathfrak{m}_n)$, déjà zariskien par le corollaire 1, est
hensélien\footnote{« Zariskien » et « hensélien » sont les mots de la page pour
le couple $(A_n, \mathfrak{p}_n)$, et ce sont les termes d'aujourd'hui. La page
ajoute que $(P^\varepsilon)'(\zeta)$, inversible dans $k$, est « en fait dans
$k_0^*$ » : il est dans $k_0$, et son inverse aussi, par le corollaire 1 pour
$n = 0$. Une première démonstration, par
l'automorphisme $(x, z) \mapsto (x, z^{N} + f_1(x) z^{N-1} + \cdots)$, a été
encadrée et biffée.}.

\textbf{Corollaire 3.} Soit $P \in A_{n+1}$, vu comme série en
$(t_1, \ldots, t_n, Z)$, avec $P(0) = 0$ et $\frac{\partial P}{\partial Z}(0)$
inversible dans $k$. L'unique $F \in S_n$ tel que $P(t, F(t)) = 0$ et
$F(0) = 0$ est dans $A_n$.

En effet, $u(x, z) = (x, P(x, z))$ est un $A$-morphisme de $X_{n+1}$ compatible
aux projections sur $X_n$ et un isomorphisme formel ; par $(\mathrm{H})$ son
inverse est de la forme $(x, z) \mapsto (x, \varphi(x, z))$ avec
$\varphi \in A_{n+1}$, et $F = \varphi(t, 0) \in A_n$. Le corollaire 3 contient le
corollaire 2 (translater $Z$ de $\zeta$) : c'est le théorème des fonctions
implicites dans la classe.

La marge de la page 12 pose une question que les feuillets laissent ouverte :
les conditions $(\mathrm{I})$, $(\mathrm{S})$, $(\mathrm{R})$, $(\mathrm{H})$
entraînent-elles que $S_n$ soit fidèlement plat sur $A_n$\footnote{La fin de
la question porte une restriction sur $k$ par rapport à $k_0$, d'une lecture
incertaine. Pour les séries algébriques sur un corps, la réponse est oui :
$S_n$ est le complété de l'anneau local noethérien $A_n$.} ?

\pagerange{13}{13}
\section*{Sous-schémas formels (page 13)}

Soit $X$ un $A$-schéma formel et $Y \subset X$ un sous-schéma formel lisse. On
dit que $Y$ est un \emph{sous-$A$-schéma formel} s'il existe sur $Y$ une
$A$-structure pour laquelle l'inclusion est un $A$-morphisme. Une telle
structure est unique. On trouve en effet un $A$-morphisme $X \to Y'$ vers un
$A$-schéma formel $Y'$ dont la restriction $Y \to Y'$ est un isomorphisme
formel\footnote{La page dit « il est immédiat ». Sur un corps, ou sur un anneau
local, une projection sur des coordonnées convenables convient ; sur un
anneau quelconque, un sous-espace tangent n'a pas toujours de supplémentaire
de coordonnées, et il faut raisonner localement sur $k$.} ; alors $Y$ est un
sous-$A$-schéma formel si et seulement si le composé
$Y' \xrightarrow{\sim} Y \hookrightarrow X$ est un $A$-morphisme, et la structure
de $Y$ est celle transportée de $Y'$ — par $(\mathrm{H})$\footnote{La page
invoque ici « (d) », d'une lecture incertaine ; c'est bien la condition
hensélienne qui sert.}.

\textbf{Exemples.} La diagonale de $X \times X$ ; plus généralement celle de
$X \times_Y X$ si $X \to Y$ est un $A$-morphisme submersif. Et, en marge : un
morphisme formel $X \to Y$ entre $A$-schémas formels est un $A$-morphisme si et
seulement si son graphe est un sous-$A$-schéma formel de $X \times Y$ (la
projection du graphe sur $X$ est un $A$-isomorphisme, par $(\mathrm{H})$).

\textbf{Schémas au-dessus de l'anneau des $A$-fonctions.} Soit $\mathcal{X}$ un
schéma sur $\mathrm{Spec}\, \mathcal{O}_A(X)$, muni d'une section $s_0$, et lisse
aux points de cette section. Le complété formel $\widehat{\mathcal{X}}$ de
$\mathcal{X} \otimes_{\mathcal{O}_A(X)} \mathcal{O}(X)$ le long de la section déduite
$s$ porte une $A$-structure pour laquelle $X \to \widehat{\mathcal{X}} \to X$ sont
des $A$-morphismes ; $\widehat{\mathcal{X}}$ dépend fonctoriellement du triple
$(X, \mathcal{X}, s_0)$, et sa formation commute aux images inverses par
$X' \to X$\footnote{La page dit « une construction utilisant e) permet d'y
mettre une $A$-structure », sans la donner, et la lettre est incertaine. La
construction naturelle choisit des coordonnées étales le long de la section
et utilise les corollaires de la condition hensélienne pour que deux choix
diffèrent d'un $A$-isomorphisme.}.

\pagerange{14}{14}
\section*{La condition sur les coefficients et le calcul différentiel (page 14)}

\begin{itemize}
  \item $(\mathrm{D})$ \emph{condition différentielle} : chaque $A_n$ est stable
    par les opérateurs différentiels élémentaires
    $D_p = \frac{1}{p_1! \cdots p_n!}\, \frac{\partial^{|p|}}{\partial t_1^{p_1}
    \cdots \partial t_n^{p_n}}$ ; de façon équivalente, si
    $f(x_1, \ldots, x_n, y_1, \ldots, y_m) \in A_{n+m}$ s'écrit
    $\sum_{p \in \mathbf{N}^m} f_p(x)\, y^p$, les coefficients $f_p$ sont dans $A_n$.
\end{itemize}

Les $D_p$ sont les dérivations de Hasse, qui ont un sens sur tout anneau.
L'équivalence est la formule de Taylor : $f(x + y) = \sum_p (D_p f)(x)\, y^p$,
où $f(x+y) \in A_{2n}$ par $(\mathrm{S})$ ; dans l'autre sens, $f_p(x) =
(D^{(y)}_p f)(x, 0)$.

La page marque la lettre de cette condition « d) », qui est déjà celle de la
condition hensélienne ; elle est vraisemblablement celle qu'invoquent les pages
13 et 15 sous la lettre e).

\textbf{Conséquences.} Pour un $A$-schéma formel $X$, les algèbres de parties
principales $\mathcal{P}^r_{X/k} = \mathcal{O}(X \times X)/\mathcal{I}_\Delta^{r+1}$ ont
une $A$-structure canonique et fonctorielle : en coordonnées, c'est le module
libre de base les $(y - x)^p$, $|p| \leq r$, sur $A_n$, et $(\mathrm{D})$ dit
précisément que l'image de $\mathcal{O}_A(X \times X)$ dans $\mathcal{P}^r_{X/k}$ est ce
module. De même pour $\mathcal{P}^r_{X/Y}$ si $X \to Y$ est un $A$-morphisme
submersif, et pour les invariants normaux $\mathcal{O}_X/\mathcal{J}^r$ d'un couple
$Y \subset X$ de $A$-schémas formels, compatiblement aux images inverses
transversales. L'opérateur différentiel universel
$\mathcal{O}(X) \to \mathcal{P}^r_{X/k}$ envoie $\mathcal{O}_A(X)$ dans
$(\mathcal{P}^r_{X/k})_A$. Les opérateurs différentiels $\mathrm{Diff}^r(X/k)$, duaux
des $\mathcal{P}^r_{X/k}$, ont des $A$-structures ; la composition respecte les
$A$-opérateurs, qui stabilisent $\mathcal{O}_A(X)$ ; de même pour les opérateurs
relatifs à un $A$-morphisme submersif\footnote{La page note l'ordre des
parties principales par l'exposant $n$, qui désigne ailleurs la dimension ;
nous écrivons $r$. Une note marginale encerclée, en haut à gauche, est pour
l'essentiel illisible.}.

\textbf{Corollaire.} $A_n \subset k_0[[t_1, \ldots, t_n]]$ : les coefficients
$(D_p f)(0)$ d'un $f \in A_n$ sont dans $A_0 = k_0$. D'où un foncteur des
$A$-schémas formels vers les schémas formels sur $k_0$ — la classe $A$ est
toujours définie sur le corps, ou l'anneau, de ses constantes.

\pagerange{15}{16}
\section*{Quotients et condition intégrale (pages 15 et 16)}

Soit $p : X \to Y$ un morphisme submersif de schémas formels lisses, $X$ muni
d'une $A$-structure. On dit que $Y$ admet une \emph{$A$-structure quotient} s'il
admet une $A$-structure faisant de $p$ un $A$-morphisme. Elle est unique : il
existe un sous-$A$-schéma formel $Y' \subset X$ (une transversale aux fibres) tel
que $Y' \to Y$ soit un isomorphisme formel ; la $A$-structure quotient existe si
et seulement si le composé $X \to Y \xrightarrow{\sim} Y'$ est un $A$-morphisme,
et c'est alors celle de $Y'$\footnote{La page invoque ici e), d'une lecture
incertaine. Le passage de $Y$ à $Y'$ et retour est un isomorphisme formel
dont on veut l'inverse dans la classe : c'est $(\mathrm{H})$ qui sert.}.

\textbf{Relations d'équivalence.} Si $p : X \to Y$ est submersif,
$R = X \mathbin{\widehat{\times}_Y} X \subset X \mathbin{\widehat{\times}} X$ est un
sous-schéma formel lisse, et une relation d'équivalence : symétrique, contenant
la diagonale, transitive. Plus précisément, $X$ étant donné, la page affirme
une correspondance biunivoque entre quotients formels submersifs $Y$ et
relations d'équivalence formelles lisses $R$\footnote{L'énoncé n'est pas
démontré. En caractéristique nulle, il se déduit du théorème de Frobenius
formel appliqué à la distribution que définit $R$.}. Si $Y$ admet une
$A$-structure quotient, $R$ est un sous-$A$-schéma formel de
$X \mathbin{\widehat{\times}} X$. On demande la réciproque :

\begin{itemize}
  \item $(\mathrm{Q})$ \emph{condition intégrale} : si $R$ est une relation
    d'équivalence formelle sur $X$ qui est un sous-$A$-schéma formel de
    $X \mathbin{\widehat{\times}} X$, le quotient formel $Y = X/R$ admet une
    $A$-structure quotient.
\end{itemize}

C'est la forme, dans la seconde rédaction, de la condition $(\mathrm{F})$ des
feuilletages de la page 6, et Grothendieck le dit en ces termes :
heuristiquement, étant donné un $A$-feuilletage de $X$, la variété formelle des
feuilles admet une $A$-structure quotient ; en caractéristique nulle, les
$A$-feuilletages correspondent aux systèmes de Pfaff complètement intégrables,
et la condition intégrale signifie que l'intégration de ce système — par le
morphisme formel $X \to Y'$ de tout à l'heure — se fait par un $A$-morphisme. La
page s'arrête sur cette phrase ; le reste du feuillet est blanc.

\pagerange{18}{19}
\section*{Le théorème de préparation formel (pages 18 et 19)}

\textbf{Théorème} (« de préparation “formel” »)\textbf{.} Soit $A$ un anneau
linéairement topologisé, séparé et complet, $Z$ une indéterminée,
$F = \sum_{i \geq 0} a_i Z^i \in A[[Z]]$, et $N \in \mathbf{N}$ tel que

\begin{itemize}
  \item[(i)] $a_i$ est topologiquement nilpotent pour $0 \leq i \leq N-1$ ;
  \item[(ii)] $a_N$ est inversible.
\end{itemize}

Alors $F$ est un élément régulier de $A[[Z]]$, et $A[[Z]]/F A[[Z]]$ est un
$A$-module libre de base les images des $Z^i$, $0 \leq i \leq N-1$ : tout
$G \in A[[Z]]$ s'écrit de façon unique
\[
  G = QF + \sum_{0 \leq i \leq N-1} b_i Z^i, \qquad Q \in A[[Z]],\ b_i \in A.
\]

C'est la division de Weierstrass sous sa forme formelle\footnote{Le mot
« division » n'est pas sur la page, qui dit « préparation ». La lettre $A$
désigne ici un anneau quelconque et non plus une classe de séries ; nous
gardons la notation de la page, qui ne rattache pas ce théorème aux pages
précédentes. Sous cette forme, sur un anneau linéairement topologisé, il
figure chez Bourbaki (\emph{Algèbre commutative}, chap.~VII, § 3), sauf erreur
de notre part sur le lieu exact.}.

\textbf{Démonstration.} $A$ est limite projective de ses quotients discrets
$A/I$, $I$ parcourant les idéaux ouverts, et $A[[Z]] = \varprojlim (A/I)[[Z]]$.
L'unicité fait que les écritures sur les $A/I$ se recollent, et un $H$ tel que
$FH = 0$ est nul modulo chaque $I$ : on peut supposer $A$ discret. Alors (i)
signifie que $a_0, \ldots, a_{N-1}$ sont nilpotents.

Soit $\mathfrak{m}$ l'idéal qu'ils engendrent : il est de type fini et engendré
par des nilpotents, donc nilpotent. Posons $A_0 = A/\mathfrak{m}$ et notons
$\dot{a}_i$, $\dot{F}$ les réductions modulo $\mathfrak{m}$. On a $\dot{a}_i = 0$
pour $i < N$, donc
\[
  \dot{F} = \dot{a}_N Z^N + \dot{a}_{N+1} Z^{N+1} + \cdots = Z^N \dot{u},
\]
où $\dot{u} = \dot{a}_N + \dot{a}_{N+1} Z + \cdots$ est une unité de $A_0[[Z]]$
puisque $\dot{a}_N$ est inversible\footnote{La page écrit
$\dot{F} = \dot{a}_N \alpha Z^N$ en mettant $\alpha$ dans $A_0^*$ ; c'est une unité
de $A_0[[Z]]$, de terme constant $1$, et non une constante. La parenthèse est
surchargée à cet endroit.}. Considérons l'application $A$-linéaire
\[
  \varphi : A^N \times A[[Z]] \longrightarrow A[[Z]], \qquad
  \bigl((b_i)_{0 \leq i \leq N-1}, Q\bigr) \longmapsto
  \sum_{0 \leq i \leq N-1} b_i Z^i + QF.
\]
Elle respecte les filtrations $\mathfrak{m}$-adiques, qui sont finies, et
$\mathfrak{m}^j A[[Z]] = \mathfrak{m}^j[[Z]]$ puisque $\mathfrak{m}$ est de type fini.
En degré $j$, l'application graduée associée est, pour le $A_0$-module
$M = \mathfrak{m}^j/\mathfrak{m}^{j+1}$,
\[
  M^N \times M[[Z]] \longrightarrow M[[Z]], \qquad
  \bigl((b_i), Q\bigr) \longmapsto \sum_{0 \leq i \leq N-1} b_i Z^i + Q\, \dot{u}\, Z^N,
\]
qui est manifestement bijective : on sépare les termes de degré $< N$ et l'on
divise le reste par $Z^N$, puis par l'unité $\dot{u}$. Une application filtrée
dont le gradué est bijectif, pour des filtrations finies, est bijective ;
$\varphi$ l'est donc. La bijectivité donne l'écriture unique, et la régularité
de $F$ : $FH = 0$ s'écrit $\varphi(0, H) = \varphi(0, 0)$, d'où $H = 0$.

Cette démonstration est la seconde de la page. Une première, encadrée et
barrée, procédait en trois temps : les $Z^i$, $i < N$, engendrent
$A[[Z]]/F A[[Z]]$ parce qu'ils l'engendrent modulo $\mathfrak{m}$ et par
Nakayama ; $F$ est régulier, par le gradué associé ; enfin, dans le cas
polynomial $F = Z^N + a_{N-1} Z^{N-1} + \cdots + a_0$, l'application
$A[Z]/F A[Z] \to A[[Z]]/F A[[Z]]$ est surjective, et bijective si et seulement
si $F A[[Z]] \cap A[Z] = F A[Z]$. L'argument par l'application $\varphi$ remplace
les trois à la fois.

\textbf{Ce que les feuillets ne disent pas.} Les pages 18 et 19 ne mentionnent
ni les classes $A$ ni les $A$-schémas formels. Le lien est pourtant direct :
demander qu'une classe de séries soit stable par la division de Weierstrass —
le quotient et le reste d'une série de la classe par une série de la classe
restant dans la classe — est l'axiome central de ce qu'on appelle depuis
Denef et Lipshitz, vers 1980, un \emph{système de Weierstrass}, et la plupart
des conditions des pages 9 à 14 (inversion des unités, fonctions implicites,
coefficients) en découlent ou l'accompagnent\footnote{Ces travaux sont
postérieurs de plus de dix ans à la datation de l'inventaire ; les feuillets
ne pouvaient pas les connaître, et nous ne suggérons pas qu'ils les
annoncent. Le rapprochement, et l'hypothèse qu'il explique la présence du
théorème dans le dossier, sont de cette lecture.}.

\end{document}
